\chapter{Interaksi Fundamental}\label{ch:g11:fundamental-interactions}

Gosokkan sebuah balon pada rambutmu dan balon itu memungut serpihan
kertas --- mengalahkan tarikan \omterm{def:g11:fundamental-interactions:four}{gravitasi} seluruh planet dengan beberapa
sentimeter persegi karet. Kemenangan yang mudah itulah sebabnya \omterm{def:g11:fundamental-interactions:ladder}{atom}
tetap utuh dan sebabnya \omterm{def:g11:fundamental-interactions:ladder}{inti atom} yang berat akhirnya berkeping
(\cref{ch:g11:nucleus-radioactivity}). Alam semesta berjalan tepat di
atas empat interaksi; bab ini menemui keempatnya, menimbangnya satu
terhadap yang lain, lalu memberi masing-masing lantainya sendiri.

\section{Materi dari kuark sampai galaksi}

\begin{definition}[Tangga struktur]\label{def:g11:fundamental-interactions:ladder}
Materi dibangun berlantai-lantai, masing-masing dirakit dari lantai di
bawahnya: tiga \emph{kuark}\index{kuark} (berupa titik, lebih kecil
daripada \qty{e-18}{m}) terikat menjadi sebuah
\emph{nukleon}\index{nukleon} --- proton atau neutron, berukuran sekitar
\qty{e-15}{m}; nukleon terkemas menjadi \emph{inti
atom}\index{inti atom} (\qty{e-15}{m} sampai \qty{e-14}{m}); inti atom
beserta awan elektronnya menjadi sebuah \emph{atom}\index{atom}
(\qty{e-10}{m}); atom terikat menjadi \emph{molekul}\index{molekul} dan
kristal (\qty{e-9}{m} ke atas), yang terakit menjadi benda sehari-hari
(\qty{1}{m}), planet (\qty{e7}{m}), sistem keplanetan (\qty{e11}{m}),
galaksi (\qty{e21}{m}) dan alam semesta teramati (\qty{e26}{m}).
\end{definition}

\begin{definition}[Keempat interaksinya]\label{def:g11:fundamental-interactions:four}
Sebuah \emph{interaksi fundamental}\index{interaksi fundamental} adalah
\omterm{def:g10:inertia:force}{gaya} yang tidak dapat dikembalikan kepada \omterm{def:g10:inertia:force}{gaya} yang lain; jumlahnya
tepat empat: \emph{gravitasi}\index{gravitasi}, \emph{interaksi
elektromagnetik}\index{interaksi elektromagnetik}, \omterm{def:g11:fundamental-interactions:strong}{interaksi kuat} dan
\omterm{def:g11:fundamental-interactions:weak}{interaksi lemah}.
\end{definition}

\begin{remark}\label{rem:g11:fundamental-interactions:empty}
Lantainya dipisahkan oleh kekosongan: \omterm{def:g11:fundamental-interactions:ladder}{atom} $10^{5}$ kali lebih lebar
daripada \omterm{def:g11:fundamental-interactions:ladder}{inti atomnya} --- besarkanlah \omterm{def:g11:fundamental-interactions:ladder}{inti atomnya} menjadi kelereng
\qty{1}{cm} dan awan elektronnya selebar satu kilometer. Sesuatu tentu
menahan lantai yang renggang itu tetap bersatu, dan hampir tidak pernah
\omterm{def:g11:fundamental-interactions:four}{gravitasi} yang melakukannya.
\end{remark}

\begin{omfigure}
\begin{tikzpicture}[font=\small, x=0.24cm]
  \draw[-Stealth] (-1,0) -- (46.5,0);
  \foreach \x/\l/\d in {0/{10^{-18}}/2pt, 3/{10^{-15}}/13pt,
                        8/{10^{-10}}/2pt, 18/{10^{0}}/2pt,
                        25/{10^{7}}/2pt, 29/{10^{11}}/13pt,
                        39/{10^{21}}/2pt, 44/{10^{26}}/13pt}
    \draw (\x,0.1) -- (\x,-0.1) node[below=\d] {$\l$};
  \foreach \x/\n in {0/kuark, 3/inti atom, 8/atom, 18/manusia, 25/Bumi,
                     29/{Matahari--Bumi}, 39/galaksi, 44/alam semesta}{
    \fill[omDef] (\x,0) circle (1.6pt);
    \node[rotate=45, anchor=west, inner sep=1.5pt] at (\x,0.24) {\n};}
\end{tikzpicture}

{\small Tangga struktur pada sumbu pangkat sepuluh: empat puluh empat
dekade memisahkan \omterm{def:g11:fundamental-interactions:ladder}{kuark} dari alam semesta teramati. Ukurannya dalam
\omterm{def:g10:orders-of-magnitude:unit}{meter}.}
\end{omfigure}

\section{Muatan listrik}

\begin{definition}[Muatan listrik]\label{def:g11:fundamental-interactions:charge}
\emph{Muatan listrik}\index{muatan listrik} adalah sifat materi yang
menjadi tempat bekerjanya \omterm{def:g11:fundamental-interactions:four}{interaksi elektromagnetik}, sebagaimana massa
adalah sifat tempat bekerjanya \omterm{def:g11:fundamental-interactions:four}{gravitasi}. Muatan diukur dalam
\emph{coulomb}\index{coulomb} (\unit{C}) dan hadir dalam \emph{dua
tanda}, positif dan negatif, yang saling menghapus ketika dijumlahkan;
muatan sejenis tolak-menolak, muatan tak sejenis tarik-menarik.
\end{definition}

\begin{definition}[Muatan dasar]\label{def:g11:fundamental-interactions:elementary}
Muatan bersifat berbutir: setiap muatan yang teramati adalah kelipatan
bulat \emph{muatan dasar}\index{muatan dasar}
$e = \qty{1.60e-19}{C}$. Proton membawa $+e$, elektron $-e$, dan neutron
$0$; muatan $q$ memuat $N = \abs{q}/e$ muatan dasar.
\end{definition}

\begin{proposition}[Kekekalan muatan]\label{prop:g11:fundamental-interactions:conservation}
Muatan seluruhnya pada sistem terasing tidak pernah berubah: muatan
tidak diciptakan dan tidak dimusnahkan, melainkan hanya dipindahkan ---
pada praktiknya oleh elektron, pembawa yang paling ringan.
\end{proposition}
\admitted

\begin{remark}\label{rem:g11:fundamental-interactions:conservation}
Belum pernah teramati satu pun pelanggarannya; alasan yang mendalam,
yaitu sebuah kesetangkupan elektromagnetisme, diberikan pada jilid
perguruan tinggi.
\end{remark}

\begin{definition}[Pemuatan dengan gesekan dan dengan sentuhan]\label{def:g11:fundamental-interactions:charging}
Menggosokkan dua isolator merenggut elektron dari yang satu ke yang
lain, sehingga tertinggal muatan berlawanan yang sama besarnya
(\emph{pemuatan dengan gesekan}\index{pemuatan dengan gesekan});
menyentuhkan benda bermuatan pada benda netral berarti berbagi muatan
(\emph{pemuatan dengan sentuhan}\index{pemuatan dengan sentuhan}). Hanya
elektron yang berpindah; muatan seluruhnya tetap kekal.
\end{definition}

\begin{example}[Menghitung elektron]\label{ex:g11:fundamental-interactions:balloon}
Balon yang digosokkan pada rambut memperoleh $q = \qty{-1.6e-8}{C}$:
balon itu \emph{memperoleh} $N = \num{1.6e-8} / \num{1.6e-19} = \num{1.0e11}$ elektron, dan rambutnya tertinggal dengan tepat
$+\qty{1.6e-8}{C}$.
\end{example}

\section{Hukum Coulomb}

\begin{theorem}[Hukum Coulomb]\label{thm:g11:fundamental-interactions:coulomb}
Dua muatan titik $q_1$ dan $q_2$ yang terpisah sejauh $d$ saling
mengerjakan \omterm{def:g10:inertia:force}{gaya} sepanjang garis penghubungnya, dengan besar yang sama
dan arah yang berlawanan --- tolak-menolak untuk tanda yang sejenis,
tarik-menarik untuk tanda yang tak sejenis --- yaitu
\[
F = k\,\frac{\abs{q_1 q_2}}{d^{2}},
\qquad k = \qty{8.99e9}{N.m^2/C^2}.
\]
\end{theorem}
\admitted

\begin{remark}\label{rem:g11:fundamental-interactions:torsion}
Coulomb menetapkan hukum ini pada tahun 1785 dengan mengukur puntiran
serat puntir yang \omterm{def:g11:lenses-and-eye:lens}{tipis}; mengapa \omterm{def:g10:orders-of-magnitude:scinot}{eksponennya} tepat $2$ dijawab pada
jilid perguruan tinggi.
\end{remark}

\begin{omfigure}
\begin{tikzpicture}[font=\small]
  \draw[omDef, thick, fill=omDef!20] (0,0) circle (0.3) node[black] {$+$};
  \draw[omDef, thick, fill=omDef!20] (2.6,0) circle (0.3) node[black] {$+$};
  \draw[-Stealth, omThm, very thick] (-0.3,0) -- (-1.5,0)
    node[midway, above] {$\vect F_{2/1}$};
  \draw[-Stealth, omThm, very thick] (2.9,0) -- (4.1,0)
    node[midway, above] {$\vect F_{1/2}$};
\end{tikzpicture}\qquad
\begin{tikzpicture}[font=\small]
  \draw[omDef, thick, fill=omDef!20] (0,0) circle (0.3) node[black] {$+$};
  \draw[omProp, thick, fill=omProp!25] (2.6,0) circle (0.3) node[black] {$-$};
  \draw[-Stealth, omThm, very thick] (0.3,0) -- (1.1,0)
    node[midway, above] {$\vect F_{2/1}$};
  \draw[-Stealth, omThm, very thick] (2.3,0) -- (1.5,0)
    node[midway, above] {$\vect F_{1/2}$};
\end{tikzpicture}

{\small Muatan sejenis tolak-menolak (kiri), muatan tak sejenis
tarik-menarik (kanan); kedua \omterm{def:g10:inertia:force}{gayanya} selalu sama besar dan berlawanan
arah.}
\end{omfigure}

\begin{remark}[Kembaran gravitasi secara bentuk]\label{rem:g11:fundamental-interactions:twin}
Hukum Coulomb, lambang demi lambang, adalah hukum \omterm{def:g11:fundamental-interactions:four}{gravitasi} semesta
(\cref{ch:g10:universal-gravitation}) dengan muatan menggantikan massa:
\begin{center}
\begin{tabular}{lll}
sumber & massa $m$ & muatan $q$ \\
hukum & $F = G\,\dfrac{m_1 m_2}{d^2}$ & $F = k\,\dfrac{\abs{q_1 q_2}}{d^2}$ \\[1.5ex]
tetapan & $G = \qty{6.67e-11}{N.m^2/kg^2}$ & $k = \qty{8.99e9}{N.m^2/C^2}$ \\
tanda & selalu tarik-menarik & tarik-menarik atau tolak-menolak \\
\end{tabular}
\end{center}
Kedua perbedaannya menggerakkan segala sesuatu di bawah ini: listrik
jauh lebih kuat, dan hanya listrik yang dapat menghapus dirinya sendiri.
\end{remark}

\begin{example}[Dua proton, kedua gayanya sekaligus]\label{ex:g11:fundamental-interactions:protons}
Dua proton ($m_p = \qty{1.67e-27}{kg}$) berada pada jarak
$d = \qty{1.0e-15}{m}$ --- bertetangga di dalam \omterm{def:g11:fundamental-interactions:ladder}{inti atom}. Tolakan
listriknya:
\[
F_E = \num{8.99e9} \times
\frac{(\num{1.60e-19})^{2}}{(\num{1.0e-15})^{2}}
\approx \qty{2.3e2}{N}
\]
--- setara \omterm{def:g10:universal-gravitation:weight}{berat} sebuah koper $23$ \omterm{def:g10:orders-of-magnitude:unit}{kilogram}, yang dipikul partikel
bermassa $10^{-27}$ \omterm{def:g10:orders-of-magnitude:unit}{kilogram}. Tarikan \omterm{def:g11:fundamental-interactions:four}{gravitasinya}:
$F_G = \num{6.67e-11} \times (\num{1.67e-27})^{2} /
(\num{1.0e-15})^{2} \approx \qty{1.9e-34}{N}$. Nisbahnya
$F_E / F_G = k e^{2} / (G m_p^{2}) \approx \num{1.2e36}$ tidak
bergantung pada jaraknya ($d^{2}$ saling menghapus): di antara partikel
dasar, \omterm{def:g11:fundamental-interactions:four}{gravitasi} $10^{36}$ kali terlalu lemah untuk berarti --- pada
jarak \emph{berapa pun}.
\end{example}

\section{Interaksi kuat dan interaksi lemah}

Contoh tadi meninggalkan sebuah skandal: \omterm{def:g11:fundamental-interactions:ladder}{inti atom} helium menahan dua
proton yang tolak-menolak dengan puluhan \omterm{def:g10:inertia:force}{newton} --- namun helium ada.

\begin{definition}[Interaksi kuat]\label{def:g11:fundamental-interactions:strong}
\emph{Interaksi kuat}\index{interaksi kuat} mengikat \omterm{def:g11:fundamental-interactions:ladder}{kuark} menjadi
\omterm{def:g11:fundamental-interactions:ladder}{nukleon} dan \omterm{def:g11:fundamental-interactions:ladder}{nukleon} menjadi \omterm{def:g11:fundamental-interactions:ladder}{inti atom}. Pada \qty{e-15}{m} interaksi itu
bersifat tarik-menarik dan jauh lebih kuat daripada tolakan listriknya
(ribuan \omterm{def:g10:inertia:force}{newton} antarnukleon); di luar beberapa $10^{-15}$ \omterm{def:g10:orders-of-magnitude:unit}{meter}
interaksi itu lenyap: \emph{jangkauan}\index{jangkauan (sebuah
interaksi)}nya sekitar \qty{e-15}{m}. Interaksi itu mencengkeram proton
dan neutron sama saja dan tidak menghiraukan muatan.
\end{definition}

\begin{omfigure}
\begin{tikzpicture}[font=\small]
  \draw[omDef, thick, fill=omDef!20] (0,0) circle (0.34) node[black] {p};
  \draw[omDef, thick, fill=omDef!20] (3.2,0) circle (0.34) node[black] {p};
  \draw[-Stealth, omThm, thick] (-0.34,0) -- (-1.7,0)
    node[midway, above] {$F_E \approx \qty{230}{N}$};
  \draw[-Stealth, omThm, thick] (3.54,0) -- (4.9,0) node[midway, above] {$F_E$};
  \draw[-Stealth, omMeth, very thick] (0.2,0.55) -- (1.45,0.55)
    node[midway, above] {kuat};
  \draw[-Stealth, omMeth, very thick] (3.0,0.55) -- (1.75,0.55)
    node[midway, above] {kuat};
  \draw[Stealth-Stealth] (0,-0.65) -- (3.2,-0.65)
    node[midway, below] {$d = \qty{1.0e-15}{m}$};
\end{tikzpicture}

{\small Tarik-menarik dan tolak-menolak di dalam \omterm{def:g11:fundamental-interactions:ladder}{inti atom}: pada
$10^{-15}$~m tarikan kuatnya (ribuan \omterm{def:g10:inertia:force}{newton}) mengalahkan tolakan
listriknya.}
\end{omfigure}

\begin{remark}[Mengapa inti atom ada --- dan mengapa yang berat menyerah]\label{rem:g11:fundamental-interactions:nuclei}
Kedua \omterm{def:g10:inertia:force}{gaya} di dalam \omterm{def:g11:fundamental-interactions:ladder}{inti atom} berskala secara berbeda. Perekat kuatnya
berjangkauan pendek: sebuah \omterm{def:g11:fundamental-interactions:ladder}{nukleon} hanya mengikat beberapa tetangganya
di dalam \qty{e-15}{m}. Tolakan listriknya berjangkauan panjang: setiap
proton mendorong \emph{setiap proton yang lain}. Ketika \omterm{def:g11:fundamental-interactions:ladder}{inti atomnya}
tumbuh, tolakannya menumpuk lebih cepat daripada perekatnya: \omterm{def:g11:fundamental-interactions:ladder}{inti atom}
yang berat memerlukan neutron tambahan (perekat tanpa muatan), dan di
luar sekitar $80$ proton bahkan neutron pun gagal --- \omterm{def:g11:fundamental-interactions:ladder}{inti atom} yang
paling berat hidup di atas waktu yang dipinjamnya
(\cref{ch:g11:nucleus-radioactivity}).
\end{remark}

\begin{definition}[Interaksi lemah]\label{def:g11:fundamental-interactions:weak}
\emph{Interaksi lemah}\index{interaksi lemah}, yang \omterm{def:g11:fundamental-interactions:strong}{jangkauannya} lebih
pendek lagi, memungkinkan sebuah neutron berubah menjadi proton (dan
sebaliknya): interaksi itulah yang menggerakkan keradioaktifan $\beta$
pada \cref{ch:g11:nucleus-radioactivity}.
\end{definition}

\section{Siapa menguasai lantai yang mana}

\begin{method}[Mengaudit sebuah skala]\label{met:g11:fundamental-interactions:audit}
Untuk menemukan interaksi yang menguasai sebuah struktur berukuran $d$,
periksalah \omterm{def:g11:fundamental-interactions:strong}{jangkauan} dan penghapusannya:
\begin{enumerate}
  \item $d \leq \qty{e-14}{m}$: \emph{\omterm{def:g11:fundamental-interactions:strong}{interaksi kuat}} yang menguasai ---
        ia mengalahkan tolakan listriknya, dan \omterm{def:g11:fundamental-interactions:four}{gravitasinya} kalah
        $10^{36}$.
  \item dari \omterm{def:g11:fundamental-interactions:ladder}{atom} sampai gunung ($\qty{e-10}{m}$ sampai $\qty{e4}{m}$):
        \omterm{def:g10:inertia:force}{gaya} kuatnya sudah di luar \omterm{def:g11:fundamental-interactions:strong}{jangkauan}; \emph{\omterm{def:g11:fundamental-interactions:four}{interaksi
        elektromagnetik}} yang menguasai --- ikatan, kelekatan, \omterm{def:g10:inertia:inventory}{gesekan},
        sentuhan.
  \item planet dan seterusnya ($d \geq \qty{e6}{m}$): materinya netral
        sampai ketelitian yang luar biasa, sehingga listriknya menghapus
        dirinya sendiri; massa hanya bertanda satu dan selalu menjumlah
        --- \emph{\omterm{def:g11:fundamental-interactions:four}{gravitasi}} yang menguasai langit.
\end{enumerate}
\omterm{def:g11:fundamental-interactions:weak}{Interaksi lemah} tidak menguasai satu lantai pun: ia tidak mengikat apa
pun, melainkan menengahi perubahan (peluruhan $\beta$).
\end{method}

\begin{omfigure}
\begin{tikzpicture}[font=\small, x=0.26cm]
  \draw[-Stealth] (-1,0) -- (41.5,0);
  \foreach \x/\l in {3/{10^{-15}}, 8/{10^{-10}}, 18/{10^{0}},
                     24/{10^{6}}, 39/{10^{21}}}
    \draw (\x,0.1) -- (\x,-0.1) node[below] {$\l$};
  \draw[omProp, fill=omProp!35] (3,0.3) rectangle (4,0.75);
  \node[anchor=west, omProp] at (4.4,0.52) {kuat};
  \draw[omDef, fill=omDef!25] (8,0.95) rectangle (22,1.4);
  \node[omDef] at (15,1.17) {elektromagnetik};
  \draw[omThm, fill=omThm!25] (24,1.6) rectangle (39,2.05);
  \node[omThm] at (31.5,1.82) {gravitasi};
\end{tikzpicture}

{\small Siapa menguasai di mana: \omterm{def:g11:fundamental-interactions:strong}{interaksi kuat} untuk \omterm{def:g11:fundamental-interactions:ladder}{inti atom},
elektromagnetik dari \omterm{def:g11:fundamental-interactions:ladder}{atom} sampai gunung, \omterm{def:g11:fundamental-interactions:four}{gravitasi} untuk planet dan
seterusnya. Ukurannya dalam \omterm{def:g10:orders-of-magnitude:unit}{meter}.}
\end{omfigure}

\section{Latihan}

\begin{exercise}[$\star$]\label{exo:g11:fundamental-interactions:1}
Berikan \omterm{def:g10:orders-of-magnitude:oom}{orde besar} tiap ukuran berikut beserta lantainya pada tangga
itu: proton \qty{1.7e-15}{m}; \omterm{def:g11:fundamental-interactions:ladder}{molekul} air \qty{2.8e-10}{m}; sebutir
pasir \qty{2e-4}{m}; Bumi \qty{1.3e7}{m}; Bima Sakti \qty{9.5e20}{m}.
\end{exercise}

\begin{exercise}[$\star$]\label{exo:g11:fundamental-interactions:2}
Sebuah balon yang digosokkan pada rambut membawa
$q = \qty{-4.8e-9}{C}$. Apakah balon itu memperoleh atau kehilangan
elektron, dan berapa banyak? Berapa muatan rambutnya sesudah itu, dan
hukum yang mana yang mengatakannya?
\end{exercise}

\begin{exercise}[$\star$]\label{exo:g11:fundamental-interactions:3}
Muatan $q_1 = \qty{+2.0e-6}{C}$ dan $q_2 = \qty{-3.0e-6}{C}$ berjarak
\qty{30}{cm}. Hitunglah \omterm{def:g10:inertia:force}{gayanya}; tarik-menarik ataukah tolak-menolak?
Bandingkan \omterm{def:g10:inertia:force}{gaya} pada $q_1$ dan \omterm{def:g10:inertia:force}{gaya} pada $q_2$.
\end{exercise}

\begin{exercise}[$\star$]\label{exo:g11:fundamental-interactions:4}
Dua muatan tarik-menarik dengan \omterm{def:g10:inertia:force}{gaya} $F_0$. Menjadi berapa \omterm{def:g10:inertia:force}{gayanya}
apabila: jaraknya digandakan; jaraknya dibagi $3$; salah satu muatannya
digandakan; kedua muatannya \emph{dan} jaraknya digandakan?
\end{exercise}

\begin{exercise}[$\star$]\label{exo:g11:fundamental-interactions:5}
Sebutkan interaksi yang bertanggung jawab atas: (a) Bulan yang
mengorbit Bumi; (b) kelekatan sebuah kristal garam; (c) kelekatan \omterm{def:g11:fundamental-interactions:ladder}{inti
atom} helium; (d) peluruhan $\beta$ karbon-14; (e) balon pada
\cref{exo:g11:fundamental-interactions:2} yang menempel pada dinding.
\end{exercise}

\begin{exercise}[$\star\star$]\label{exo:g11:fundamental-interactions:6}
Kedua proton pada \omterm{def:g11:fundamental-interactions:ladder}{inti atom} helium berjarak $d = \qty{2.0e-15}{m}$
($m_p = \qty{1.67e-27}{kg}$). Hitunglah tolakan listriknya, tarikan
\omterm{def:g11:fundamental-interactions:four}{gravitasinya}, dan nisbahnya. Apa yang menjaga \omterm{def:g11:fundamental-interactions:ladder}{inti atom} itu tetap utuh?
\end{exercise}

\begin{exercise}[$\star\star$]\label{exo:g11:fundamental-interactions:7}
Di dalam \omterm{def:g11:fundamental-interactions:ladder}{atom} hidrogen, elektronnya ($m_e = \qty{9.11e-31}{kg}$)
berjarak $d = \qty{5.3e-11}{m}$ dari protonnya
($m_p = \qty{1.67e-27}{kg}$). Hitunglah \omterm{def:g10:inertia:force}{gaya} listrik dan \omterm{def:g10:universal-gravitation:force}{gaya gravitasi}
antara keduanya beserta nisbahnya. Apa yang menahan \omterm{def:g11:fundamental-interactions:ladder}{atom} tetap utuh?
\end{exercise}

\begin{exercise}[$\star\star$]\label{exo:g11:fundamental-interactions:8}
Bola A membawa $q_A = \qty{+8.0e-9}{C}$; bola B yang serupa bersifat
netral. Keduanya disentuhkan, lalu dipisahkan sejauh \qty{10}{cm}.
Berapa muatan yang dibawa masing-masing (periksalah kekekalan
muatannya)? Hitunglah \omterm{def:g10:inertia:force}{gaya} di antara keduanya --- tarik-menarik ataukah
tolak-menolak?
\end{exercise}

\begin{exercise}[$\star\star$]\label{exo:g11:fundamental-interactions:9}
Dua muatan serupa $q = \qty{1.0e-6}{C}$ tolak-menolak dengan
\qty{0.90}{N}: berapa jarak keduanya? Pada jarak berapa \omterm{def:g10:inertia:force}{gayanya} turun
menjadi \qty{0.10}{N}?
\end{exercise}

\begin{exercise}[$\star\star$]\label{exo:g11:fundamental-interactions:10}
Dua proton pada \omterm{def:g11:fundamental-interactions:ladder}{molekul} yang bertetangga berjarak
$d = \qty{1.0e-10}{m}$. Hitunglah tolakan listriknya. Berapa sumbangan
\omterm{def:g11:fundamental-interactions:strong}{interaksi kuat} pada jarak ini? Interaksi yang mana yang mengikat
\omterm{def:g11:fundamental-interactions:ladder}{molekul}?
\end{exercise}

\begin{exercise}[$\star\star$]\label{exo:g11:fundamental-interactions:11}
Tiap keping dari dua keping uang logam memuat sekitar $\num{e23}$
elektron. Pindahkanlah satu elektron dari tiap \emph{sejuta} elektron di
antara keduanya lalu tempatkan kedua keping itu \qty{1.0}{m} terpisah.
Hitunglah muatan tiap kepingnya, lalu \omterm{def:g10:inertia:force}{gayanya}. \omterm{def:g10:universal-gravitation:weight}{Berat} massa berapa yang
setara dengan \omterm{def:g10:inertia:force}{gaya} itu? Simpulkan tentang kenetralan materi
sehari-hari.
\end{exercise}

\begin{exercise}[$\star\star\star$]\label{exo:g11:fundamental-interactions:12}
Sebuah \omterm{def:g11:fundamental-interactions:ladder}{inti atom} uranium berdiameter \qty{1.4e-14}{m} dan berisi $92$
proton. Hitunglah tolakan antara dua proton pada ujung yang
berseberangan. Dapatkah \omterm{def:g11:fundamental-interactions:strong}{interaksi kuat} mengikat pasangan itu --- mengapa
tidak? Ada berapa pasang proton yang saling menolak? Mengapa \omterm{def:g11:fundamental-interactions:ladder}{inti atom}
yang berat memerlukan neutron tambahan, dan mengapa di luar ukuran
tertentu bahkan neutron pun tidak sanggup menyelamatkannya?
\end{exercise}

\begin{exercise}[$\star\star\star$]\label{exo:g11:fundamental-interactions:13}
Dua bola bermassa \qty{1.0}{g} tergantung pada benang \qty{50}{cm} yang
terikat pada titik yang sama. Setelah diberi muatan $q$ yang sama,
keduanya berhenti pada jarak \qty{6.0}{cm}. Tiap bola berada dalam
kesetimbangan di bawah \omterm{def:g10:universal-gravitation:weight}{beratnya}, \omterm{def:g11:circuits-and-power:voltage}{tegangan} benangnya dan \omterm{def:g10:inertia:force}{gaya}
listriknya: tunjukkan bahwa $F = mg\tan\theta$ ($\theta$: sudut benang
terhadap arah tegak) dan bahwa di sini $\tan\theta \approx 0.060$.
Hitunglah $F$, lalu $q$, lalu banyaknya \omterm{def:g11:fundamental-interactions:elementary}{muatan dasar} yang berpindah.
\end{exercise}

\begin{exercise}[$\star\star\star$]\label{exo:g11:fundamental-interactions:14}
Data: $M_E = \qty{5.97e24}{kg}$, $M_M = \qty{7.35e22}{kg}$, jarak
Bumi--Bulan $d = \qty{3.84e8}{m}$. Hitunglah \omterm{def:g10:universal-gravitation:force}{gaya gravitasi}
Bumi--Bulan. \omterm{def:g11:fundamental-interactions:four}{Gravitasinya} dimatikan: muatan $+Q$ pada Bumi dan $-Q$
pada Bulan yang berapa yang mempertahankan tarikan yang sama? Setara
berapa elektron $Q$ itu, dan berapa \omterm{def:g10:universal-gravitation:weight}{beratnya}
($m_e = \qty{9.11e-31}{kg}$)? Berikan ulasanmu.
\end{exercise}

\begin{exercise}[$\star\star\star$]\label{exo:g11:fundamental-interactions:15}
Di dalam kristal garam, ion Na$^{+}$ dan Cl$^{-}$ (muatannya $\pm e$)
berjarak $d = \qty{2.8e-10}{m}$; sebuah ion Cl$^{-}$ bermassa
\qty{5.9e-26}{kg}. Hitunglah \omterm{def:g10:inertia:force}{gaya} antara ion yang bertetangga dan \omterm{def:g10:universal-gravitation:weight}{berat}
sebuah ion Cl$^{-}$; berikan nisbahnya. Lalu jelaskan bagaimana
\omterm{def:g11:fundamental-interactions:four}{gravitasi} menang pada skala yang besar --- tidak ada gunung yang
melampaui \qty{e4}{m} --- padahal tiap ikatannya mengalahkan \omterm{def:g11:fundamental-interactions:four}{gravitasi}
sebesar lima belas \omterm{def:g10:orders-of-magnitude:oom}{orde besar}.
\end{exercise}

\section{Soal: Audit keempat gaya}

\begin{problem}[{Soal akhir pekan --- mengapa dunia tidak runtuh dan
tidak pula berhamburan: keempat interaksinya diaudit lantai demi
lantai, sampai ke kenetralan materi yang mencapai dua bagian dari
$10^{18}$}]
\label{pb:g11:fundamental-interactions:1}
\omterm{def:g11:fundamental-interactions:ladder}{Atom} tidak runtuh ke dalam, \omterm{def:g11:fundamental-interactions:ladder}{inti atom} sebagian besar bertahan, dan
Bulan tidak menabrak dan tidak pula lepas; soal ini mengaudit siapa
yang pantas menerima pujiannya. Datanya:
$e = \qty{1.60e-19}{C}$,
$k = \qty{8.99e9}{N.m^2/C^2}$, $G = \qty{6.67e-11}{N.m^2/kg^2}$,
$m_p = \qty{1.67e-27}{kg}$, $m_e = \qty{9.11e-31}{kg}$,
$M_E = \qty{5.97e24}{kg}$, $M_M = \qty{7.35e22}{kg}$, jarak
Bumi--Bulan \qty{3.84e8}{m}.

\medskip
\noindent\textbf{Bagian I --- Tangganya.}
\begin{enumerate}
  \item Daftarkan lantai tangganya dari \omterm{def:g11:fundamental-interactions:ladder}{kuark} sampai galaksi, satu \omterm{def:g10:orders-of-magnitude:oom}{orde
        besar} ukuran untuk masing-masing.
  \item Model berskala: \omterm{def:g11:fundamental-interactions:ladder}{inti atomnya} menjadi kelereng \qty{1}{cm};
        selebar apa \omterm{def:g11:fundamental-interactions:ladder}{atomnya}?
  \item Berapa bagian volume sebuah \omterm{def:g11:fundamental-interactions:ladder}{atom} yang ditempati \omterm{def:g11:fundamental-interactions:ladder}{inti atomnya}?
  \item Sehelai rambut tebalnya \qty{70}{\micro m}. Berapa \omterm{def:g11:fundamental-interactions:ladder}{atom} yang
        membentang menyeberanginya?
  \item Berapa pangkat sepuluh dari batas atas \omterm{def:g11:fundamental-interactions:ladder}{kuark} (\qty{e-18}{m})
        sampai sebuah galaksi (\qty{e21}{m})?
\end{enumerate}

\medskip
\noindent\textbf{Bagian II --- Lantai listriknya.}
\begin{enumerate}[resume]
  \item Di dalam \omterm{def:g11:fundamental-interactions:ladder}{atom} hidrogen, elektronnya berjarak
        $d = \qty{5.3e-11}{m}$ dari protonnya: hitunglah \omterm{def:g10:inertia:force}{gaya} listrik
        di antara keduanya.
  \item Hitunglah \omterm{def:g10:universal-gravitation:force}{gaya gravitasi} di antara keduanya, dan nisbah kedua
        \omterm{def:g10:inertia:force}{gayanya}. Yang mana yang menahan \omterm{def:g11:fundamental-interactions:ladder}{atomnya} tetap utuh?
  \item Apa yang berubah jika \omterm{def:g11:fundamental-interactions:four}{gravitasi} dimatikan di dalam \omterm{def:g11:fundamental-interactions:ladder}{atom}? Dan
        jika listriknya yang dimatikan?
  \item Mengapa \omterm{def:g10:inertia:force}{gaya} ``sentuh'' dalam kehidupan sehari-hari --- lantai
        yang mendorong telapak kakimu, \omterm{def:g10:inertia:inventory}{gesekan} --- sebenarnya
        elektromagnetik yang menyamar?
  \item Tiap keping dari dua keping uang logam memuat sekitar
        $\num{e23}$ elektron. Pindahkanlah satu elektron dari tiap
        \emph{semiliar} elektron di antara keduanya, lalu tempatkan
        keduanya \qty{1.0}{m} terpisah: hitunglah muatannya dan
        \omterm{def:g10:inertia:force}{gayanya}. Apa kesimpulanmu?
  \item Interaksi yang mana yang menguasai lantai dari \qty{e-10}{m}
        sampai \qty{e4}{m}, dan mengapa \omterm{def:g11:fundamental-interactions:strong}{interaksi kuat} tidak ikut
        bersaing?
\end{enumerate}

\medskip
\noindent\textbf{Bagian III --- Lantai intinya.}
\begin{enumerate}[resume]
  \item Dua proton berjarak \qty{1.0e-15}{m} di dalam sebuah \omterm{def:g11:fundamental-interactions:ladder}{inti atom}:
        hitunglah tolakan listriknya.
  \item Jika tidak tertahan, percepatan berapa yang diberikan \omterm{def:g10:inertia:force}{gaya} itu
        kepada sebuah proton? Bandingkan dengan $g$.
  \item Daftarkan tiga sifat yang diperlukan \omterm{def:g11:fundamental-interactions:strong}{interaksi kuat} untuk
        menjelaskan mengapa \omterm{def:g11:fundamental-interactions:ladder}{inti atom} ada, mengapa neutron ikut
        tercengkeram, dan mengapa \omterm{def:g11:fundamental-interactions:ladder}{inti atom} berukuran mungil.
  \item Di dalam \omterm{def:g11:fundamental-interactions:ladder}{inti atom} uranium (diameternya \qty{1.4e-14}{m}, $92$
        proton), hitunglah tolakan antara dua proton pada ujung yang
        berseberangan. Dapatkah \omterm{def:g11:fundamental-interactions:strong}{interaksi kuat} mengikat pasangan itu
        secara langsung? Jelaskan mengapa tolakannya menang ketika \omterm{def:g11:fundamental-interactions:ladder}{inti
        atomnya} tumbuh, dan untuk apa neutronnya.
  \item Dalam satu kalimat, sebagaimana dijanjikan: apa yang dikerjakan
        \omterm{def:g11:fundamental-interactions:weak}{interaksi lemah}?
\end{enumerate}

\medskip
\noindent\textbf{Bagian IV --- Lantai astronomi, dan putusannya.}
\begin{enumerate}[resume]
  \item Hitunglah \omterm{def:g10:universal-gravitation:force}{gaya gravitasi} antara Bumi dan Bulan.
  \item Bumi memuat sekitar $\num{1.8e51}$ proton (dan elektron
        sebanyak itu pula), Bulan sekitar $\num{2.2e49}$. Berapa bagian
        $f$ dari muatan proton tiap benda, yang dibiarkan tak terhapus,
        yang akan membuat \omterm{def:g10:inertia:force}{gaya} listriknya menyamai \omterm{def:g10:universal-gravitation:force}{gaya gravitasinya}?
  \item Mengapa \omterm{def:g11:fundamental-interactions:four}{gravitasi}, si lemah pada pertanyaan 7, menguasai
        astronomi? Berikan dua alasan.
  \item Putusannya --- satu kalimat untuk tiap lantai (\omterm{def:g11:fundamental-interactions:ladder}{inti atom}, \omterm{def:g11:fundamental-interactions:ladder}{atom}
        sampai gunung, planet dan seterusnya): apa yang menahannya tetap
        bersatu? Lalu jawablah judulnya: mengapa dunia tidak runtuh dan
        tidak pula berhamburan?
\end{enumerate}
\end{problem}
